% Options for packages loaded elsewhere
\PassOptionsToPackage{unicode}{hyperref}
\PassOptionsToPackage{hyphens}{url}
\PassOptionsToPackage{dvipsnames,svgnames,x11names}{xcolor}
\documentclass[
]{ltjsarticle}
\usepackage{xcolor}
\usepackage[margin=20mm]{geometry}
\usepackage{amsmath,amssymb}
\setcounter{secnumdepth}{-\maxdimen} % remove section numbering
\usepackage{iftex}
\ifPDFTeX
  \usepackage[T1]{fontenc}
  \usepackage[utf8]{inputenc}
  \usepackage{textcomp} % provide euro and other symbols
\else % if luatex or xetex
  \usepackage{unicode-math} % this also loads fontspec
  \defaultfontfeatures{Scale=MatchLowercase}
  \defaultfontfeatures[\rmfamily]{Ligatures=TeX,Scale=1}
\fi
\usepackage{lmodern}
\ifPDFTeX\else
  % xetex/luatex font selection
\fi
% Use upquote if available, for straight quotes in verbatim environments
\IfFileExists{upquote.sty}{\usepackage{upquote}}{}
\IfFileExists{microtype.sty}{% use microtype if available
  \usepackage[]{microtype}
  \UseMicrotypeSet[protrusion]{basicmath} % disable protrusion for tt fonts
}{}
\makeatletter
\@ifundefined{KOMAClassName}{% if non-KOMA class
  \IfFileExists{parskip.sty}{%
    \usepackage{parskip}
  }{% else
    \setlength{\parindent}{0pt}
    \setlength{\parskip}{6pt plus 2pt minus 1pt}}
}{% if KOMA class
  \KOMAoptions{parskip=half}}
\makeatother
\usepackage{longtable,booktabs,array}
\newcounter{none} % for unnumbered tables
\usepackage{calc} % for calculating minipage widths
% Correct order of tables after \paragraph or \subparagraph
\usepackage{etoolbox}
\makeatletter
\patchcmd\longtable{\par}{\if@noskipsec\mbox{}\fi\par}{}{}
\makeatother
% Allow footnotes in longtable head/foot
\IfFileExists{footnotehyper.sty}{\usepackage{footnotehyper}}{\usepackage{footnote}}
\makesavenoteenv{longtable}
\setlength{\emergencystretch}{3em} % prevent overfull lines
\providecommand{\tightlist}{%
  \setlength{\itemsep}{0pt}\setlength{\parskip}{0pt}}
\usepackage{bookmark}
\IfFileExists{xurl.sty}{\usepackage{xurl}}{} % add URL line breaks if available
\urlstyle{same}
\hypersetup{
  colorlinks=true,
  linkcolor={Maroon},
  filecolor={Maroon},
  citecolor={Blue},
  urlcolor={Blue},
  pdfcreator={LaTeX via pandoc}}

\author{}
\date{}

\begin{document}

\section{論文16:測度問題の厳密な縮約 ---
干渉粒度・イベント構造・条件付けへの有限分解}\label{ux8ad6ux658716ux6e2cux5ea6ux554fux984cux306eux53b3ux5bc6ux306aux7e2eux7d04-ux5e72ux6e09ux7c92ux5ea6ux30a4ux30d9ux30f3ux30c8ux69cbux9020ux6761ux4ef6ux4ed8ux3051ux3078ux306eux6709ux9650ux5206ux89e3}

\textbf{著者}: 木原範昭 \textbf{作成}: 2026年6月12日 \textbf{版}:
v1.1(v1.0 に §12
観察追記{[}集計一意化問題と集合選択問題の対応候補の発見・検討{]}＋外部参考文献を追加。経緯の透明性のための追記であり、§1.2
の方針・本文の結果は不変。v1.0=第2回査読{[}採録・軽微修正条件付き{]}R1〜R6
反映、v0.2=第1回査読{[}中改訂{]}反映) \textbf{DOI(本版)}:
10.5281/zenodo.20686057 / \textbf{Concept DOI}: 10.5281/zenodo.20665699
\textbf{ライセンス}: CC BY 4.0 \textbf{シリーズ}:
波長空間と周波数空間の双対幾何(論文1〜15 の続編、第16篇)

\begin{center}\rule{0.5\linewidth}{0.5pt}\end{center}

\subsection{要旨}\label{ux8981ux65e8}

体系は経路振幅(輸送位相 η、論文14)を自由パラメータなしで与える ---
確率の\textbf{素材}は完全に揃っている。決まっていないのは素材をどう束ねるかの\textbf{集計規則}であり、これは標準理論が
Born
則として公理で置く当のものである。本論文は集計規則を\textbf{選定しない}。本論文の主結果は、選定問題そのものの\textbf{厳密な縮約}である:(1)
\textbf{歴史的な近接 0.0007 は意味のある一致ではなかった}(定理1):s=9
の双子観測量は、排他による2結果強制と両測度の支配結果集中の\textbf{二重圧縮}で全変動距離が潰れる構造であり、3結果空間が開く
s=11 で乖離は 0.24 へ\textbf{3桁跳躍}する。(2)
\textbf{構造的乖離}(定理2):導出測度(振幅)は数え上げ測度(配置数)に古典化しない
--- 乖離は s とともに恒常的に O(0.1〜0.9) であり、集計の規約幅(graded
変種)でも消えない。\textbf{両測度は台(ゼロの位置)を共有し、重みでのみ乖離する}(3シェル双方向)。(3)
\textbf{干渉粒度は規約に還元されない}(定理3・非規約性):「どの代替をコヒーレントに足すか」の選択
--- チャネル一貫(歴史的 W)か配置粒度(W′)か ---
は規約ではない。鏡像親の厳密例:\((-3,0,0,0)\)
は\textbf{チャネル一貫では完全に消え(W=0
厳密)、配置粒度では完全に正常(W′=768=W′(A)
厳密)}。しかもこの相殺はスケール自由でなく(m=2
で不成立)、その発生条件は\textbf{四値ロック}(\(z^-_K/z^+_K\in\{0,\infty,\pm1\}\)、36配置全数・例外ゼロ)という\textbf{証明対象の代数}に還元済みである。(4)
\textbf{条件付けは枝チャネル層を要求する}(定理4):双子測度は親対の関係クラスの関数であり(全クラスで一定、例外ゼロ)、強度縞が読めない相対符号(dot
の符号)に依存する --- 読み出し三層
11⊂13⊂21(論文13)のうち\textbf{層3(枝チャネル)だけが測度を支える}。(5)
\textbf{イベント構造はちょうど1ビット}:一回(合併キー)か逐次(分割キー)かは記録原理を介して干渉粒度に翻訳され、数え上げ側の歴史的二分(一括/逐次)はその位相なし影である
--- 2×2表が数値で完成している。(6)
\textbf{総括(主定理)}:集計規則の残る自由度は\textbf{干渉粒度 ×
イベント構造1ビット ×
枝チャネル条件付け}の有限な決定構造に縮約され、候補間の判定は\textbf{モデル内部で計算可能な実験}(判定表
§9)として立っている。付録に、この線で証明済みの相殺定理群(η\(^{2}\)
等分布の対合 σ**・A1 シェル対合・B2
トーサー残差×因数分解・鏡像反転補題と指標和還元)を、監査済みの範囲限定つきで収める。Born
則は本縮約の下で「然るべき極限で合流すべき有効則」の地位に置かれ、合格基準は\textbf{収束・抑制(統計的検閲
\textasciitilde1/√N)・予言}の三つ組である。物理的同一視は行わない。

\begin{center}\rule{0.5\linewidth}{0.5pt}\end{center}

\subsection{1. 序論}\label{ux5e8fux8ad6}

\subsubsection{1.1
問いと、本論文が選定をしない理由}\label{ux554fux3044ux3068ux672cux8ad6ux6587ux304cux9078ux5b9aux3092ux3057ux306aux3044ux7406ux7531}

論文5〜15
で、状態・保存・時間・位置・位相・座標は定理または機械検証として確立した。残るのは:振幅
\(\eta\) の集合から\textbf{確率}を作る規則 ---
どの代替をコヒーレントに足し(粒度)、何で条件付け、どのイベント構造で読むか
--- である。標準理論はここを Born
則という\textbf{公理}で済ませる。本体系は逆に、候補となる集計が\textbf{導出できすぎる}:逐次・一括(数え上げ)、チャネル一貫・配置粒度(振幅)、一回・逐次(結合振幅)---
いずれも在庫から定義でき、有限の s
で\textbf{真に異なる値}を与える。したがって現段階の誠実な成果は「どれが正しいか」の宣言ではなく、\textbf{選定問題を有限の決定構造に縮約し、判定実験を内部に立てる}ことである。それが本論文である。

\subsubsection{1.2
主張しないこと}\label{ux4e3bux5f35ux3057ux306aux3044ux3053ux3068}

\begin{enumerate}
\def\labelenumi{(\roman{enumi})}
\tightlist
\item
  集計規則の選定(本論文の全結果は選定に依存しない)。(ii) Born
  則の導出(地位の整理まで --- §10)。(iii)
  偏差の実在(構造の所在の特定まで)。(iv) 物理的同一視。
\end{enumerate}

\subsection{2. 定義(一般 s
の測度族)}\label{ux5b9aux7fa9ux4e00ux822c-s-ux306eux6e2cux5ea6ux65cf}

双子 \((s,s)\) に対し(容量
\(c(m)=|SH[m]|\)、チャネル=\texttt{all\_finals(s)}):

\begin{itemize}
\tightlist
\item
  \textbf{D1(逐次数え上げ)}: 先手チャネル確率 ∝
  残容量上の配置数、後手は条件付き更新。良定義性の鍵は補題「条件付き配置数は殻ごとの占有数のみに依存する」(容量二項係数の積)。
\item
  \textbf{D2(一括数え上げ)}: 結合終配置の等重。
\item
  \textbf{D3(導出双子)}:
  \(P\propto \mathrm{mult}\cdot W_{F_1}W_{F_2}\cdot[\)結合容量整合\(]\)。\(W_F=|z_F|^2\)
  は正準規約(論文14 定理5)のチャネル一貫振幅。
\item
  \textbf{W′(配置粒度)}: \(W'_F=\sum_K|z_K|^2\)(終配置 \(K\)
  ごとに二乗してから和)。
\item
  \textbf{D4(追跡量)}:
  \(\Delta(s)=\mathrm{TV}(\text{D3},\text{D1})\)、\(\Delta'(s)\) は W′
  版。
\end{itemize}

アンカー(s=9)全点厳密再現:逐次 396/403、一括 60/61、導出
0.98195、\(\Delta(9)=0.00068\)、単一崩壊 24/31。

\subsection{3. 定理1:0.0007
の正体は二結果圧縮}\label{ux5b9aux740610.0007-ux306eux6b63ux4f53ux306fux4e8cux7d50ux679cux5727ux7e2e}

\textbf{圧縮補題(一行)}: 二つの確率分布が同一の支配結果に質量
\(\ge1-\varepsilon\) を置くとき
\(\mathrm{TV}=1-\sum\min(P,Q)\le\varepsilon\)(締めた形は査読者による)。s=9
の双子では (i) 排他(原点の唯一性)が結果空間を2結果に強制し、(ii)
両測度とも支配結果に \textasciitilde0.98
集中(\(\varepsilon\approx0.018\))--- この二重圧縮が TV
を潰していた(図1a)。s=11 の3結果空間で圧縮が解けると
\(\Delta(11)=0.24136\) ---
\textbf{3桁の跳躍}(図1b)。傍証:単一崩壊水準では s=9 でも 0.774 vs 0.965
と大差;セクターB は s=9 でも \(\Delta_B=0.030\)。\textbf{0.98195 と
0.98263 の近接は s=9・セクターA
限定の偶然}であり、歴史的な「測度候補に近接」の意義づけはここで訂正される(値は不変)。

\subsection{4.
定理2:構造的乖離(古典化しない)}\label{ux5b9aux74062ux69cbux9020ux7684ux4e56ux96e2ux53e4ux5178ux5316ux3057ux306aux3044}

\textbf{検証域の明示}:本定理は \(s\in\{9,11,13\}\)
の全数水準(3点系列)の言明である。\(\Delta(s)\) はこの系列で 0.0007 →
0.24 → 最大
0.90(親型・セクター依存が大;図2)。数え上げ側内部の幅(TV(逐次,一括))は常に一桁小さい。容量整合を比率化する規約変種(graded)でも乖離
O(0.1) は消えない --- \textbf{乖離は規約幅に対して頑健}。独立の一般 s
観測量 \(D(s)=\mathrm{TV}\)(コヒーレント, インコヒーレント) も
0.135→0.140→0.185
と縮まらない。一方、\textbf{台の一致}:容量不能チャネルと振幅零チャネルは
s=9/11/13 で双方向に一致((1,1,s−2) 型のみ)---
両測度は\textbf{ゼロの位置を共有し、重みでのみ乖離する}。

\subsection{5.
定理3:干渉粒度の非規約性}\label{ux5b9aux74063ux5e72ux6e09ux7c92ux5ea6ux306eux975eux898fux7d04ux6027}

\textbf{内部定義(規約)}:
本シリーズで「規約」とは、簿記の値を変えない表示の選択をいう(例:対角・分岐の扱い
--- 論文14 定理5 の粒度監査)。チャネル一貫和(歴史的
W:記録で区別可能な終配置を跨いでコヒーレント)と配置粒度(W′:配置ごとに二乗)の選択は、この意味の規約では\textbf{ない}
--- 値を代数的に分岐させる(図3a):

{\def\LTcaptype{none} % do not increment counter
\begin{longtable}[]{@{}lll@{}}
\toprule\noalign{}
親 & W(チャネル一貫) & W′(配置粒度) \\
\midrule\noalign{}
\endhead
\bottomrule\noalign{}
\endlastfoot
\((+3,0,0,0)\)@13 & 9216 & 768(30配置) \\
\((-3,0,0,0)\)@13(鏡像) & \textbf{0(厳密)} & \textbf{768 = W′(A)
厳密} \\
\end{longtable}
}

鏡像親は配置粒度では完全に正常に崩壊し、チャネル一貫の集計だけが健全な分布を厳密零まで歪める。さらに(補遺84
の m=4 実行を §5 本文に反映 --- R1):(i)
この相殺は\textbf{スケール自由でない}(m=2 では W(B)=32≠0、W′
も不一致)--- 発生条件が m に依存する。ただし
\textbf{m=4(s=21)で相殺の単位がチャネルであることが判明}:鏡像は全体としては消えない(\(S_+-S_-\neq0\)、パリティ仮説支持)が、\textbf{チャネル
(5,7,9) は単体で相殺}(\(W_A=147456\)、\(W_B=0\)、\(W'\) は両者 3072)---
偶 m でもチャネル単位の \(W=0\wedge W'>0\)
が出現し、粒度間乖離の供給源は奇 m 部分列より\textbf{豊富}である。また
\textbf{W′ の鏡像同値は m≥3 で全チャネル厳密一致}(m=2 の不一致は小 s
の例外と見られる)。(ii) その m 依存性は\textbf{鏡像反転補題}
\(\eta_B/\eta_A=-i(-1)^{\sigma_0}\)(全数
394経路・反例ゼロ)により\textbf{σパリティ指標和の消滅} \(S_+=S_-\)
に還元され、配置水準では\textbf{四値ロック}
\(z^-_K/z^+_K\in\{0,\infty,+1,-1\}\)(m=3
全36配置:12+12+6+6、例外ゼロ;排他族の重み対
384=384;図3b)という証明対象の代数に固定されている。(iii)
セクター(ホロノミー)依存は配置粒度でほぼ消える(\(W'\) は A≈B)---
\textbf{ホロノミーの情報は配置間の相対位相に局在する}。

\textbf{用語注記}: 本論文の鏡像対の「A/B」(親 \(\pm m\))と、論文14 §5
のホロノミー類「A/B」は\textbf{別物}である --- 併読時の混同に注意。

\subsection{6.
定理4:条件付けは枝チャネル層を要求する}\label{ux5b9aux74064ux6761ux4ef6ux4ed8ux3051ux306fux679dux30c1ux30e3ux30cdux30ebux5c64ux3092ux8981ux6c42ux3059ux308b}

双子の結合振幅は親対の相対配置に強く依存する(完全相殺対を含む)。検査の三層区別(査読
\#4):\textbf{クラス内一定(s=9 セクターA
23対・全9クラス・例外ゼロ)は①構成により真}である --- 不変量クラスは
B\(_{4}\)×swap
軌道の厳密な分類(補遺76/77)であり、一定性は計算機械の同変性の帰結で反証力を持たない(整合性確認)。\textbf{落ちえた検査(③)＝本節の核心}はこちらである:クラスのうち
dot=±1 の対は\textbf{強度縞では区別できないのに P(X) が異なる}(0.972 vs
0.994 --- 等しいこともできた)。よって測度は読み出し三層 11⊂13⊂21(論文13
定理3)の\textbf{層1〜2
では支えられず、層3(枝チャネル)を要求する}。「干渉粒度=記録区別能」の原理は、層を明示してのみ成立する。床(軸台が交わらない対=配置粒度値
0.96644)も③:粒度監査と独立に一致した。

\subsection{7.
イベント構造:ちょうど1ビット}\label{ux30a4ux30d9ux30f3ux30c8ux69cbux9020ux3061ux3087ux3046ux30691ux30d3ux30c3ux30c8}

一回イベント(中間追記なし→合併占有集合→分割横断コヒーレント)か、逐次(第一追記が分割を確定)か
--- 記録原理がこの1ビットを干渉粒度に翻訳する。数え上げの歴史的二分(一括
60/61・逐次 396/403)はその位相なし影であり、全体は 2×2
表に収まる(数値完成):

{\def\LTcaptype{none} % do not increment counter
\begin{longtable}[]{@{}
  >{\raggedright\arraybackslash}p{(\linewidth - 4\tabcolsep) * \real{0.3333}}
  >{\raggedright\arraybackslash}p{(\linewidth - 4\tabcolsep) * \real{0.3333}}
  >{\raggedright\arraybackslash}p{(\linewidth - 4\tabcolsep) * \real{0.3333}}@{}}
\toprule\noalign{}
\begin{minipage}[b]{\linewidth}\raggedright
\end{minipage} & \begin{minipage}[b]{\linewidth}\raggedright
一回
\end{minipage} & \begin{minipage}[b]{\linewidth}\raggedright
逐次
\end{minipage} \\
\midrule\noalign{}
\endhead
\bottomrule\noalign{}
\endlastfoot
数え上げ & 0.98361 & 0.98263 \\
振幅 & 対依存:平均 0.989、幅 {[}0.966, 1.000{]}+完全相殺対(両キー残存) &
対依存:平均 0.781、幅 {[}0, 1{]} \\
\end{longtable}
}

振幅側はどちらのイベント構造でも対配置依存が残る ---
条件付け(§6)が分離不能の成分であることの表現。

\subsection{8.
証明済みの相殺定理群(範囲限定つき)}\label{ux8a3cux660eux6e08ux307fux306eux76f8ux6bbaux5b9aux7406ux7fa4ux7bc4ux56f2ux9650ux5b9aux3064ux304d}

この線の探究で、チャネル一貫和の上の厳密相殺が\textbf{対合・指標和の様式}で証明された:(i)
\(\Sigma\eta^2=0\)(区分的対合 σ** ---
完全性・固定点なし・\(\eta\to\pm i\eta\);\textbf{s=9 全数水準})。(ii)
A1(\(\Sigma i^{\delta_1}=0\))--- シェル水準の対合
\(x\leftrightarrow x\mp2\)(借入↔貸出)。(iii) B2(sin枝 \(Z_4\) 等分布)---
トーサー残差(4解軌道=完全代表系)×段の因数分解;\textbf{531
経路集合のファイバー構造に基づき、333 終は等分布しない}。(iv)
鏡像反転補題と指標和還元(§5)。これらは集計規則の候補(チャネル一貫和)の\textbf{内部構造}であり、その測度論的地位の裁定(本縮約)と相補的である。

\subsection{9.
主定理(候補族の分類)と判定表}\label{ux4e3bux5b9aux7406ux5019ux88dcux65cfux306eux5206ux985eux3068ux5224ux5b9aux8868}

\begin{quote}
\textbf{定理6(候補族の分類 --- 本論文の主定理)}:
在庫から定義された集計候補の族 --- \{数え上げ D1/D2、振幅のチャネル一貫
W、配置粒度 W′\}×\{一回、逐次\}×\{関係クラス条件付けの有無\} ---
について:(i) 族の任意の二候補は有限 \(s\)(\(s\le13\)
の全数水準)で総変動距離 \(O(0.1)\)
の差を持ち\textbf{得る}(真に異なる測度である)。(ii)
全候補は台の一致(§4)と正準規約(論文14 定理5)を共有する。(iii)
\textbf{独立に定義された全候補がこの直積に重複なく整理され、三軸の各々は独立である}
--- どの軸も単独で変えると測度が変わる実例を持つ(粒度:鏡像親
W=0∧W′\textgreater0/イベント構造:0.98361 vs 0.98263・結合振幅の平均
0.989 vs 0.781/条件付け:dot=±1
の不等)。本定理は\textbf{族内の分類定理}であり、「集計候補はこれに尽きる」という族外の全称は主張しない(査読
\#1 による限定;(iii) の非自明な内容は重複なき整理と軸の独立性 --- R2)。
\end{quote}

\subsubsection{判定表(モデル内部で計算可能な実験)}\label{ux5224ux5b9aux8868ux30e2ux30c7ux30ebux5185ux90e8ux3067ux8a08ux7b97ux53efux80fdux306aux5b9fux9a13}

{\def\LTcaptype{none} % do not increment counter
\begin{longtable}[]{@{}
  >{\raggedright\arraybackslash}p{(\linewidth - 6\tabcolsep) * \real{0.2500}}
  >{\raggedright\arraybackslash}p{(\linewidth - 6\tabcolsep) * \real{0.2500}}
  >{\raggedright\arraybackslash}p{(\linewidth - 6\tabcolsep) * \real{0.2500}}
  >{\raggedright\arraybackslash}p{(\linewidth - 6\tabcolsep) * \real{0.2500}}@{}}
\toprule\noalign{}
\begin{minipage}[b]{\linewidth}\raggedright
\#
\end{minipage} & \begin{minipage}[b]{\linewidth}\raggedright
判定実験
\end{minipage} & \begin{minipage}[b]{\linewidth}\raggedright
何が決まるか
\end{minipage} & \begin{minipage}[b]{\linewidth}\raggedright
状態
\end{minipage} \\
\midrule\noalign{}
\endhead
\bottomrule\noalign{}
\endlastfoot
1 & パリティ仮説:m 奇 \(\Longleftrightarrow\) \(S_+=S_-\) &
鏡像相殺の発生条件 & \textbf{m=4 実行済み(補遺84):支持} ---
\(S_+-S_-\neq0\)、ただしチャネル (5,7,9) は偶 m
でも相殺(単位はチャネル別指標和に細分)。\textbf{W′ の鏡像同値は m≥3
で一般的}(全チャネル厳密一致)。m=5@s31 が残(Born 負荷点) \\
2 & 候補測度の大 s 相互収束(\(\Delta/\Delta'\) 系列の延長、粒度別) &
測度選択が「深い離散領域のみ」で効くか「全域」で効くか & 未実行 \\
3 & ι 構成(隠れ vx の軸0枝反転) & 四値ロックの証明部品 &
\textbf{実行済み(補遺84):m=3 で完全成立} --- 配置保存・σ反転
384/384・乗数 ±1(192/192)。\textbf{m=2 では像の集合閉性が 2/10 に破れる}
---「なぜ m 偶で壊れるか」の構造的住所=ι
の閉性。残:乗数の配置内分布則 \\
4 & ゲージ独立性(マーキング・向き付けに対する測度の不変性) & 原理 v2
の反証可能内容①(クラス内一定は同変性で自動のため) & 未実行 \\
5 & s=11 非飽和域の同型対比較 & dot=±2 等値が飽和の偶然か追加対称性か &
未実行 \\
\end{longtable}
}

\subsection{10. Born
則の地位と合格基準}\label{born-ux5247ux306eux5730ux4f4dux3068ux5408ux683cux57faux6e96}

本縮約の下で、Born
則は「再現すべき厳密な的」ではなく\textbf{極限と検閲精度の中で合流すべき有効則}である。合格基準は三つ組:\textbf{収束}(選定測度が然るべき極限で
Born に合流 ---
極限写像自体が構成課題)、\textbf{抑制}(検証済み領域で偏差が\textbf{統計的検閲}
\textasciitilde1/√N 以下 --- 論文15
の検閲量子の統計版)、\textbf{予言}(検閲を超える偏差が住む領域=深い離散領域での検証可能な構造
---
粒度・パリティ・量子化)。下限は確保されている:内部原理で決まらなければ集計規則を一つ公準として置けばよく、それで公理勘定は標準理論と同水準に並ぶ(下回らない)。上振れは、記録原理(枝チャネル層)が集計を一意化する場合であり、そのとき
Born の類似物は定理になる。

\subsection{11. 正直な限界}\label{ux6b63ux76f4ux306aux9650ux754c}

\begin{enumerate}
\def\labelenumi{\arabic{enumi}.}
\tightlist
\item
  \textbf{選定はしていない}(§1.2)---
  本論文の縮約は選定の前提整理であり、決定構造の3成分のどれが先に固定されるかは判定表の実行に依る。
\item
  数値系列は s≤13(一部 s=11)の全数事実。大 s の挙動は判定表2。
\item
  結合振幅の親対掃引は s=9・セクターA(23対)。s≥11
  の掃引・混合セクター対は未実行。
\item
  四値ロック・重み対は m=3 の全数事実(証明は判定表3の仕様で進行中)。
\item
  物理的同一視は行わない。
\end{enumerate}

\subsection{12.
観察追記(v1.1):集計一意化問題と集合選択問題の対応候補}\label{ux89b3ux5bdfux8ffdux8a18v1.1ux96c6ux8a08ux4e00ux610fux5316ux554fux984cux3068ux96c6ux5408ux9078ux629eux554fux984cux306eux5bfeux5fdcux5019ux88dc}

v1.0(前版)の時点では気づいていなかったが、その後に先行研究を発見して内容を検討したところ、§10
で Born
接続として残した未解決の問い------「記録原理(枝チャネル層)が集計規則を一意化するか」------が、量子力学基礎論で半世紀にわたり扱われてきた\textbf{集合選択問題(set
selection
problem)}と同種の課題に起因している可能性に気づいた。経緯の透明性のため、この発見と検討の事実を観察として追記する(同定の正否は現時点で保留)。

\begin{itemize}
\tightlist
\item
  §10 の「記録原理が集計を一意化すれば Born
  の類似物は定理になる」は、構造として、記録・準古典性が好ましい歴史集合を選ぶという
  Gell-Mann--Hartle の構想 {[}2,3{]}
  と同型に見える。対応先は整合的/デコヒーレント歴史の枠組み {[}1,4{]}
  と、その物理的基盤であるデコヒーレンス・einselection {[}5{]} である。
\item
  その一般的障害は、整合性条件が確率を割り当てる基準を与えるものの一意の歴史集合を選ばないという\textbf{集合選択問題}
  {[}6,7{]}
  にあたる可能性がある。これが正しければ、本論文の集計一意化は本論文に固有の困難ではなく、約30年未解決の既知の問題に対応していることになる。
\item
  一方、本論文がこの非一意性を有限の決定構造(干渉粒度 ×
  イベント構造1ビット ×
  条件付け、§5--9)に縮約し、候補間の判別をモデル内部の判定実験として立てている点は、抽象的非一意性で止まる一般論と異なって見える。これが実質的前進か離散・有限設定の産物かは、本論文では判定しない。
\end{itemize}

\textbf{留保}:この同定は未検証の観察であり、表層的類似にすぎない可能性、対応のどこかが破れている可能性を排除しない。また記録原理は標準量子力学でも必ずしも成り立たない例が知られる(witness
が存在しても情報転写なしに干渉が消える例
{[}8{]})。詳細な対応表・限界・全参照は補遺85(GitHub)に記録した。本追記は
§1.2 の方針(集計規則の選定をしない/Born
則を導出しない)を一切変更せず、本文の全結果も不変である。物理的同一視は行わない。

\subsection{図}\label{ux56f3}

\begin{itemize}
\tightlist
\item
  \textbf{図1}(\texttt{paper16\_fig1\_compression.png}): 二結果圧縮 ---
  s=9(2結果・Δ=0.00068)と s=11(3結果・Δ=0.24136)の結合分布
\item
  \textbf{図2}(\texttt{paper16\_fig2\_divergence.png}): 乖離の地形 ---
  Δ(s) 全親型系列と配置粒度 Δ′(セクター依存の縮退)
\item
  \textbf{図3}(\texttt{paper16\_fig3\_granularity.png}): 粒度の非規約性
  --- 鏡像相殺の m 系列(m=2 不成立/m=3 全体/m=4 はチャネル (5,7,9)
  単体)と W′ 鏡像同値(m≥3);四値ロック 12+12+6+6 と重み対 384=384
\end{itemize}

全図は機械検証済み数値の提示(模式図なし)。

\subsection{付録A:検証一覧}\label{ux4ed8ux9332aux691cux8a3cux4e00ux89a7}

{\def\LTcaptype{none} % do not increment counter
\begin{longtable}[]{@{}
  >{\raggedright\arraybackslash}p{(\linewidth - 6\tabcolsep) * \real{0.2500}}
  >{\raggedright\arraybackslash}p{(\linewidth - 6\tabcolsep) * \real{0.2500}}
  >{\raggedright\arraybackslash}p{(\linewidth - 6\tabcolsep) * \real{0.2500}}
  >{\raggedright\arraybackslash}p{(\linewidth - 6\tabcolsep) * \real{0.2500}}@{}}
\toprule\noalign{}
\begin{minipage}[b]{\linewidth}\raggedright
主張
\end{minipage} & \begin{minipage}[b]{\linewidth}\raggedright
ラベル(①構成により真/②定理+確認/③落ちえた検査)
\end{minipage} & \begin{minipage}[b]{\linewidth}\raggedright
検証
\end{minipage} & \begin{minipage}[b]{\linewidth}\raggedright
方式
\end{minipage} \\
\midrule\noalign{}
\endhead
\bottomrule\noalign{}
\endlastfoot
アンカー再現 & ③(再現失敗があり得た) & 全点一致 & 厳密(分数演算) \\
Δ(9/11/13) 全親型 & ③ & 表(図2) & 全数 \\
台の一致 & ③ & s=9/11/13 双方向 & 全数 \\
graded 規約幅 & ③ & s=9/11/13 & 全数 \\
鏡像親 W=0∧W′=768 / m=2,4 & ③ & m=1〜4 & 全数 \\
クラス一定性 & \textbf{①(同変性の帰結 --- 整合性確認)} & 23対×全9クラス
& 全数 \\
縞クラステスト(層3要求) & \textbf{③(本節の核心)} & dot=±1 不等 & 全数 \\
2×2表(A\_seq 含む) & ③ & 完成 & 全数/掃引 \\
相殺定理群(σ**・A1・B2) & ②(対合構成=証明) & s=9 水準 & 全数+対合 \\
鏡像反転補題・指標和還元 & ②(3行算術=証明) & 394経路 反例0 & 全数 \\
四値ロック・重み対 & ③(現状; ι で②化が進行) & 36配置 例外0・384=384 &
全数 \\
\end{longtable}
}

\begin{center}\rule{0.5\linewidth}{0.5pt}\end{center}

\subsection{参考文献(外部)}\label{ux53c2ux8003ux6587ux732eux5916ux90e8}

本シリーズは外部文献を引用しない方針だが、本追記(§12)は先行研究の発見経緯を誠実に記録する性格上、隣接領域の最小限の外部参照を持つ(概念的背景であり、本論文の縮約結果の根拠ではない)。全文献は補遺85
を参照。

\begin{itemize}
\tightlist
\item
  {[}1{]} R. B. Griffiths, ``Consistent histories and the interpretation
  of quantum mechanics,'' \emph{J. Stat. Phys.} \textbf{36}, 219 (1984).
\item
  {[}2{]} M. Gell-Mann, J. B. Hartle, ``Quantum mechanics in the light
  of quantum cosmology,'' in \emph{Complexity, Entropy, and the Physics
  of Information}, ed.~W. H. Zurek (Addison-Wesley, 1990).
\item
  {[}3{]} M. Gell-Mann, J. B. Hartle, ``Classical equations for quantum
  systems,'' \emph{Phys. Rev.~D} \textbf{47}, 3345 (1993).
\item
  {[}4{]} R. Omnès, \emph{The Interpretation of Quantum Mechanics}
  (Princeton Univ. Press, 1994).
\item
  {[}5{]} W. H. Zurek, ``Decoherence, einselection, and the quantum
  origins of the classical,'' \emph{Rev.~Mod. Phys.} \textbf{75}, 715
  (2003).
\item
  {[}6{]} F. Dowker, A. Kent, ``On the consistent histories approach to
  quantum mechanics,'' \emph{J. Stat. Phys.} \textbf{82}, 1575 (1996).
\item
  {[}7{]} A. Kent, ``Consistent sets yield contrary inferences in
  quantum theory,'' \emph{Phys. Rev.~Lett.} \textbf{78}, 2874 (1997).
\item
  {[}8{]} C. S. Lent, ``Blind witnesses quench quantum interference
  without transfer of which-path information,'' \emph{Entropy}
  \textbf{22}, 776 (2020).
\end{itemize}

\begin{center}\rule{0.5\linewidth}{0.5pt}\end{center}

\textbf{謝辞・履歴}: 検証は Claude Code(ローカル機械検証)と
claude.ai(独立検算・複数命題の提案と相互訂正)の二者独立検算プロトコルによる。素材:補遺68〜84(2026-06-12;補遺83=四値ロック検証記録、補遺84=m=4/ι
実行)。v0.2:第1回査読 \#1〜\#6
反映。v1.0:第2回査読(採録・軽微条件付き)R1〜R6 反映 --- 補遺84
結果の本文・図3への伝播、定理6(iii)
の非自明化、圧縮補題、分布幅、用語注記。本論文の探究では予言の反証(スケール自由性)・説明の訂正(σ\(_{0}\)定数性)・定義の照合(縞三層)が二者間で計7回往復した
--- 縮約の信頼性はこの相互訂正の履歴に依拠する。

物理的同一視は行わない。

\end{document}
